% !TeX root = ../papers/nc-p3.tex
\section{Walsh表示の既約三次元族と指定を忘れた一般商}
\label{sec:walsh-presentation-space}
\begin{theorem}\label{thm:walsh-presentation-space}
三色ずつを持つWalsh窓の欠色をともに3にして射影化する。
その表示空間は既約三次元で、幾何学的一般点は反対階数16、最小周期2の
標準Koszul左右AS正則窓である。全ての正則なWalsh窓は最小周期が2を割る。
次の指定なし一般商も次元3である。$N_{be}=(-1)^{(b+e)_0+(b+e)_1}m_{be}^2$ とし
\[
 \textstyle \sum_b N_{be}=\sum_eN_{be}=0,\qquad
 P^2=\prod_{b\ne e}N_{be},\qquad (N,P)\sim(\lambda N,\lambda^3 P).
\]
この二重被覆を行列添字 $0,1,2$ の同時置換 $S_3$ で割ったものが、
一般の添字保持同型類を与える。移動は $(N,P)\mapsto (N^{\mathsf t},-P)$ である。
正則な階数12,14,15窓はこの幾何学的一般点からの特殊化であり、
いずれも全次数で平坦な付値環上の族で実現できる。
\end{theorem}
\begin{proof}
行和・列和が零の $3\times3$ 行列の射影空間は $\mathbb P^3$ である。
九つの成分は相異なる非比例な線型式である。第一成分を基準にした他の八成分の比は
平方類として独立である。各成分が消える超平面の付値で当該指数の偶奇を読めばよい。
射影座標ごとの平方写像は有限平坦であり、逆像の一般ファイバーは次数256の体である。
平坦性によりその閉包が全逆像であり、三色ずつの開条件を取っても既約である。
座標を三つの対角成分と $N_{01}-N_{10}$ に取れる。後者が非零なら
反対階数は $12$ に対角三成分のFourier変換の階数を足したものである。
対角の平方が独立な一般点では四つの文字和は非零である。

この既約族は前二節の正則点を含む。各一文字切断が全階数で隣接階数が6の近傍で、
テンソルとその縮約から双対経路の長さ3,4の次元には $4,1$ の下界がある。
正則点での値と上半連続性から近傍で各々一定であり、長さ5は零となる。
そこで得た有限平坦なFrobenius双対の内部次数別にbar複体を取ると、
定理\ref{thm:reversible-component}の議論で幾何学的一般点のKoszul性と左右AS性が従う。
この段階では全次数について単一の開集合を主張する必要はない。
一般点と任意の上記正則点で各次数のHilbert関数が等しいので、正則点を中心とし
関数体を分数体とする離散付値環（例えばその点の吹上げ・正規化の例外素因子）では
各次数の有限生成経路加群が自由である。

全ラベルへの平行移動 $x$（偶位置に $(-1)^{x_1i_0}$ を補う）と全位置での
文字倍が位数16の射影群 $G$ を与える。位数は一位置の作用だけでも16である。
一般の階数15点の射影安定化群は定理\ref{thm:rank-fifteen-pauli-comparison}により
ちょうど $G$ である。この際、反対核の非退化な二行列で残りの二基底変更を
射影的に決めるので、独立共役群での安定化群は四位置でのものに等しい。
幾何学的一般点でも安定化群はちょうど $G$ であることを説明する。

まず複素数上で四基底を正規直交とする。各一位置のGram行列はラベルの異なる
二切片の台が交わらないため対角で、平行移動が係数の絶対値を保つのでスカラーである。
従って全ての表示は $\mathrm{SL}_4^4$ のノルム最小ベクトルである。実際各因子の
跡零Hermite行列による一径数変換で、ノルム平方は実指数関数の非負和で、
その微分は原点で零である。射影安定化群が有限なら全安定化元は射影ユニタリである。
極分解の正値部分に沿うノルムの最小値が二点で一致すれば一定であり、
正値部分の一径数群がベクトルを固定するからである（最小ノルムの等しい軌道上で
スカラー倍率の絶対値は1）。安定化Lie環の零性は上の階数15点の近傍で開である。
そこへ余分な安定化元を持つ点が解析的に集積するなら、コンパクトな射影ユニタリ群で
部分列は階数15点の安定化元へ収束する。その $G$ の元を除き、単位元へ収束する
非自明な安定化元を得る。単位元近傍の対数をその長さで割り部分列を取れば
極限の非零なLie環元が階数15点の直線を固定し矛盾する。
余分な安定化元の存在は有限型の軌道方程式から構成可能であるため、
この解析的近傍の存在でザリスキー一般点でも排除できる。全条件は代数的なので標数零へ移る。

従って一般の二表示の同型は固定した $G$ を正規化する。$G$ の元を $(x,\chi)$ とすると、
偶頂点と奇頂点での行列の交換子形式の差は
$x_1y_0+y_1x_0$ である。その根基 $x=0$ は通常の対角文字群で、
四本の同時固有線が各頂点で定まる。従って一般の同型は単項行列である。
全九成分が非零の台を保つ置換を数える。四点の全置換は $E$ のアフィン変換で、
$0,1,2$ の全ての $b,e$ について四文字の和が零となる条件から、
四位置の線型部分は同じ $A\in\mathrm{GL}_2(\mathbb F_2)$ である。
欠色3を保つと位置0,2で同じ平行移動、位置1,3でも同じ平行移動となる。
行と列の色には同じ置換 $b\mapsto Ab+3+A3$ が働き、六置換全てが出る。

この置換の符号を補うには双線型形式
$T(x,y)=x_1y_0+(Ax)_1(Ay)_0$ を使う。これは対称であり、
$q(x+y)/(q(x)q(y))=(-1)^{T(x,y)}$ となる四乗根値の関数 $q$ がある。
$q(0)=1$ とし二つの基底で $q(e_j)^2=(-1)^{T(e_j,e_j)}$ を選べばよい。
$i,k=i+d$ の基底倍率で符号を補うと、係数に残るのは $q(d)$ またはその逆と $d$ の文字だけである。
$q(d)^2=s_d/s_{Ad}$ なので $N$ はそのまま置換される。
非対角六成分では各 $d\ne0$ が二回出て $\prod_{d\ne0}q(d)^2=1$ となり、$P=\prod_{b\ne e}m_{be}$ も保つ。

位置別の対角変換の効果を決める。ある $(b,e)$ の四語で基底倍率の積を取ると
各位置の全倍率を一度ずつ使うため、全係数比の四乗は一定である。
共通倍率を除いて平方の比は符号となる。三つの非零の行成分の和が零で、
符号変更後も零なら、三符号は同一である。列にも同じ議論を適用して
平方の共通倍率を除き九係数の符号だけが残る。
36語の各四位置で用いた基底を記録する整数重み行列のSmith因子は
1が十一個、2が二個である。その整数左核の基底を各色対の符号へ制限すると、
非零となる条件は六つの非対角符号の積が1という一式だけである。
従って $N,P$ は共通倍率を除いて必要十分であり、六置換を加えれば一般の指定なし同型が決まる。
\path{computations/verify_walsh_presentation_space.py} は重み行列と符号補正を再現する。

一位置移動は $m_{be}\mapsto(-1)^{b_1(b+e)_0}m_{eb}$ であり、
六つの非対角成分で符号の積は $-1$ である。一般の $N$ はその転置の同時置換・定数倍ではない
（対角三値を一般に取れば置換と倍率を固定し、$N_{01}-N_{10}\ne0$ とする）。
二位置移動の符号 $(-1)^{d_1d_0}$ は、全位置の文字 $i$ を
$(-1)^{i_0i_1}$ 倍すると消える。実際 $b,e\ne3$ より $b_0 b_1=e_0 e_1=0$ である。
この基底変換は各係数が零であっても有効であり、全正則点で周期は2を割る。
\end{proof}

\begin{proposition}\label{prop:walsh-dual-square-cut}
三色ずつを持つ任意のWalsh候補で、その二次双対は、元の窓に各位置の
四つの欠色語 $x_{p,i}x_{p+1,i+3}$ を零とする関係を追加した商と、
矢の対角基底変更で一致する。
\end{proposition}
\begin{proof}
$D=\operatorname{diag}(1,-1,-1,1)$ と置く。隣接ブロックを
$M=\sum z_dP_d$ とすれば $M^{\mathsf t}DM=MDM^{\mathsf t}=(\sum s_d z_d^2)D$ である。
各非零ブロックは階数2で右辺は零なので、関係平面 $R$ に対し $R^\perp=DR$ となる。
行ブロック、列ブロック、転置ブロックとも同じ恒等式で処理できる。
欠色では $R=0$ のため双対関係は四語全体である。
$D$ は偶位置の矢だけの文字倍で全切断に同時に実現する。
右の文字に掛かる場合の余分な $s_b$ は色ごとのスカラーである。
\end{proof}

\begin{theorem}\label{thm:walsh-dual-frobenius}
三色ずつを持つ全てのWalsh候補の二次双対 $B$ は各始点でHilbert関数
$(1,4,6,4,1)$ のFrobenius窓である。特にこの枝の標準AS性に残る条件はKoszul性だけである。
$B$ で各矢 $x_{p,i}$ の右零化イデアルは $x_{p+1,i+3}$ の生成する右イデアルに等しく、
左側でも前位置の $i+3$ が零化イデアルを生成する。
\end{theorem}
\begin{proof}
二次双対の三次の双対空間は元の隣接関係を左右から挿入した空間の交叉である。
三文字の和 $h$ を固定すると、語は隣接二色 $b,e$ で一意に定まる。
交叉の元は $b,e\ne3$ に台を持つ。各行では全四成分が所定の関係平面に入り、
さらに欠色の一成分が零となる。この平面は座標超平面に含まれない
（非退化三次元対角形式の全等方部分空間の次元は高々1である）。
従ってその切断は直線で、実際に最後の矢を $h$ に固定した $W$ の縮約が各行でこの直線を張る。
列でも同様である。この既知の縮約の非零成分上で各行倍率と各列倍率が一致する。
台の二部グラフは各行・各列に少なくとも二辺を持つ。これは三非零色と
$N$ の和の消滅から従う。各側が三頂点なのでグラフは連結であり、全倍率は共通となる。
巡回したテンソルでも符号と転置だけが変わり、各始点で三次は各色一次元となる。

四文字の和を $h$ とすれば、任意の四次双対の双対テンソルは
\[
        f(i)W_{i+h,j,k,t}=g(t)W_{i,j,k,t+h}
\]
の両表示を持つ。$h\ne0$ で $i+j=3$ とすれば右辺は零である。
左辺の行 $3+h$ は非零色なので、任意の $i$ で $k,t$ を選んで係数を非零にでき、
$f(i)=0$ となる。$h=0$ では $t=i+e$ で各 $e=0,1,2$ を選んで
$f(i)=g(i+e)$ を得るので $f,g$ は共通定数である。従って四次は $W$ の直線に双対である。
五次に全色 $h$ を固定すると、最初の四文字に $W$ を持つ表示は最後の文字が $h$ に固定され、
後の四文字に $W$ を持つ表示は最初の文字が $h$ に固定される。
$W$ は一つの最初の文字に台を持たないため交叉は零である。
全ての完全対合は、$W$ の単文字切断の階数4、隣接切断の階数6と上記の次元から従う。

命題\ref{prop:walsh-dual-square-cut}の表示で $x_{p,i}x_{p+1,i+3}=0$ である。
他の三文字との二次積はそれぞれ異なる色で非零である。実際関係平面は全等方であり、
一つの座標語は非等方なのでその平面に入らない。よって各単文字の左倍、右倍は
次数0,1でそれぞれ階数1,3である。四次の完全対合を使えば次数2,3の階数はそれぞれ
3,1となる。ここで相手側の写像は四文字テンソルを巡回した位置における単文字の乗法である。
隣り合う欠色の文字の乗法の合成は零であり、全次数で一方の像と他方の核の次元が一致する。
これで左右の零化イデアルの記述を得る。
\end{proof}
